\documentclass[11pt]{article}
\usepackage[a4paper,margin=26mm]{geometry}
\usepackage{amsmath,amssymb,amsthm,booktabs,array}
\usepackage[T1]{fontenc}
\usepackage{lmodern}
\usepackage[colorlinks=true,allcolors=blue,breaklinks=true]{hyperref}
\newtheorem{theorem}{Theorem}
\newtheorem{corollary}[theorem]{Corollary}
\newcommand{\Q}{\mathbb Q}
\newcommand{\xt}{\mathbf x}
\title{An alternative proof of Kurkov's logarithmic-polynomial recursion
and its Leonov--Shiryaev interpretation}
\author{Haobo Ma\\\small ChronoAI Pte Ltd, Singapore
\and Wenlin Zhang\\\small National University of Singapore\\\small ChronoAI Pte Ltd, Singapore}
\date{14 November 2026}
\begin{document}
\maketitle
\begin{abstract}
We give an alternative proof of Mikhail Kurkov's conjecture in OEIS A263634, FORMULA of
30 March 2025 (entry revision \#166, 5 November 2025), identifying its
all-column array as a specialization of the Leonov--Shiryaev
product-cumulant theorem (1959).
The quotient method appears in xpaul (2019), and IV\_ (2020)
connects its first column to logarithmic polynomials.
For Kurkov's quadratic recursive array we establish the all-column identity
$R(n,k)=(n-1)![t^{n-1}]F^{(k)}(t)/F(t)$, where
$F(t)=1+\sum_{i\geq1}x_it^i/i!$.
Differentiating these quotients gives the recurrence directly.
We also rederive classical coefficient and Appell identities, correct a
factorial normalization in the source's Appell comments, and recover
the exponential-polynomial exception of Craven and Csordas (1994).
\end{abstract}

\section{The conjecture and the quotient identity}
Kurkov asks whether the recursively defined $R(n,1)$ equals the $n$th
logarithmic polynomial. The quotient formula in Theorem~\ref{thm:main}
answers this for every $n$ and describes every positive column as well.
Work in $A=\Q[x_1,x_2,\ldots]$ and $A[[t]]$. Define $L_0=0$ and
\begin{equation}\label{eq:log}
 F(t)=1+\sum_{i\geq1}x_i\frac{t^i}{i!},\qquad
 \log F(t)=\sum_{n\geq1}L_n\frac{t^n}{n!}.
\end{equation}
The logarithm is the formal sum $\sum_{m\geq1}(-1)^{m-1}(F-1)^m/m$;
each coefficient is a finite sum. The $L_n$ are also the cumulant
polynomials recorded in OEIS A127671, in a different monomial order
\cite{oeis,oeis127671}. Kurkov's conjecture concerns a particular
quadratic recursion for these already known polynomials, not their
existence or their classical coefficient formula.
The differential recurrence behind the proof has a prior quotient
representation in xpaul's 21 March 2019 answer~\cite{xpaul};
IV\_ related its repeated derivatives to logarithmic polynomials
on 15 July 2020~\cite{iv}. We give the explicit coefficient translation
to Kurkov's array. More generally, the all-column array is a specialization
of the Leonov--Shiryaev product-cumulant theorem (1959)~\cite{ls}.
We make that specialization explicit below, using the two-factor formula
proved by Beuman, Turner and Vitelli~\cite[Appendix A, (A10)--(A15)]{btv}
and Burchardt's algebraic formulation~\cite[Section 1.2, Theorem 1.9]{burchardt}.
The same conjecture was posed as the middle case of a three-weight
question on MathOverflow on 21 March 2025~\cite{mo}; we treat the
logarithmic-polynomial weight here.
The moment recurrence is classical
\cite[Section 1.2, p.~4]{novak}, and the Appell recurrence below
was given by Copeland~\cite[Part II, p.~3]{copeland}; the Appell
generating-function framework also appears in
Anshelevich~\cite[Section 1, (1); Section 2.2]{anshelevich}.

For $n,k\geq1$, set $R(1,k)=x_k$ and, for $n\geq2$, set
\begin{equation}\label{eq:rec}
 R(n,k)=R(n-1,k+1)-\sum_{j=1}^{n-1}
 \binom{n-2}{j-1}R(j,k)R(n-j,1).
\end{equation}
These formulas define all columns by induction on $n$.

\begin{theorem}\label{thm:main}
For every $n,k\geq1$,
\begin{equation}\label{eq:quotient}
 R(n,k)=(n-1)![t^{n-1}]\frac{F^{(k)}(t)}{F(t)}.
\end{equation}
In particular $R(n,1)=L_n$. The identity remains valid after
substitution of the $x_i$ by arbitrary elements of any commutative
$\Q$-algebra.
\end{theorem}
\begin{proof}
The series $F$ is invertible because $F(0)=1$. Put $Q_k=F^{(k)}/F$.
The formal product rule, applied to $FQ_k=F^{(k)}$, gives
\begin{equation}\label{eq:diff}
 Q_k'=Q_{k+1}-Q_kQ_1,\qquad Q_k(0)=x_k\quad(k\geq1).
\end{equation}
Write $Q_k=\sum_{m\geq0}q_{m,k}t^m/m!$. For $m\geq0$,
coefficient extraction from \eqref{eq:diff} gives
\[
 q_{m+1,k}=q_{m,k+1}-\sum_{r=0}^m\binom mr q_{r,k}q_{m-r,1}.
\]
With $m=n-2$ and $j=r+1$, this is precisely \eqref{eq:rec}
for $q_{n-1,k}$, with the same initial values. Simultaneous induction
on $n$ over all $k\geq1$ proves \eqref{eq:quotient}.
Termwise differentiation of the formal logarithm gives
$(\log F)'=F'/F=Q_1$. Its coefficient of $t^{n-1}/(n-1)!$
is $L_n$, which proves the conjecture. All operations commute with
homomorphisms of commutative $\Q$-algebras, proving the last assertion.
\end{proof}

For example,
\[
 R(2,k)=x_{k+1}-x_kx_1,\qquad
 R(3,k)=x_{k+2}-2x_1x_{k+1}+2x_1^2x_k-x_2x_k.
\]
Thus $L_1=x_1$, $L_2=x_2-x_1^2$, and
$L_3=x_3-3x_1x_2+2x_1^3$, agreeing with the source examples.

\paragraph{Classical mixed-cumulant interpretation.}
Put $x_0=1$ and define the unital $A$-linear map
$\Phi:A[Z]\longrightarrow A$ by $\Phi(Z^i)=x_i$.
Joint cumulants $\kappa(Y_1,\ldots,Y_r)$ are the coefficients
$[s_1\cdots s_r]\log\Phi(\exp(\sum_i s_iY_i))$.
This is the commutative-algebra setting of
Burchardt~\cite[Section 1.2, (3)]{burchardt}; no positivity assumption
on $\Phi$ is required. Formal differentiation gives
\[
 \sum_{m\geq0}\kappa(Z^k,\underbrace{Z,\ldots,Z}_{m})\frac{t^m}{m!}
 =\left.\frac{\partial}{\partial s}\log\Phi(e^{sZ^k+tZ})\right|_{s=0}
 =\frac{F^{(k)}(t)}{F(t)}.
\]
Thus $K(n,k)=\kappa(Z^k,\underbrace{Z,\ldots,Z}_{n-1})$
equals $R(n,k)$ by Theorem~\ref{thm:main}.
The two-factor product-cumulant identity~\cite[(A10)]{btv} is
\[
 \kappa(UV,Z_1,\ldots,Z_q)
 =\kappa(U,V,Z_1,\ldots,Z_q)
 +\sum_{S\subseteq\{1,\ldots,q\}}
   \kappa(U,Z_S)\kappa(V,Z_{S^c}),
\]
where $Z_S$ denotes the list indexed by $S$, including the empty list.
Choose $U=Z^k$, $V=Z$, $q=n-2$ and $Z_i=Z$.
The left side becomes $K(n-1,k+1)$ and the first term is $K(n,k)$.
There are $\binom{n-2}{j-1}$ subsets of size $j-1$, each contributing
$K(j,k)K(n-j,1)$. Rearranging gives exactly \eqref{eq:rec}.
This identifies the whole indexed array as a specialization of the
previously proved product-cumulant formula.

\section{Coefficient and Appell consequences}
\begin{corollary}\label{cor:coeff}
If $a_1,\ldots,a_n\geq0$, $\sum_{i=1}^n ia_i=n$, and
$m=\sum_{i=1}^na_i$, then
\begin{equation}\label{eq:coeff}
 [x_1^{a_1}\cdots x_n^{a_n}]L_n
 =\frac{(-1)^{m-1}(m-1)!\,n!}{\prod_{i=1}^n a_i!(i!)^{a_i}}.
\end{equation}
Every partition of $n$ contributes exactly one nonzero monomial.
Consequently $L_n$ has $p(n)$ terms, and $L_n(1,\ldots,1)=0$
for $n>1$, with $L_1(1)=1$.
\end{corollary}
\begin{proof}
In the defining sum for $\log F$, this monomial occurs only in
$(F-1)^m$. Its coefficient there is $m!/\prod a_i!$ times
$\prod(i!)^{-a_i}$. Multiplying by $(-1)^{m-1}/m$ and by $n!$
proves \eqref{eq:coeff}. The displayed number is nonzero. The exponent
vectors are exactly the integer partitions of $n$ in multiplicity
notation. Substituting $x_i=1$ in \eqref{eq:log} gives $F=e^t$ and
$\log F=t$, proving the row-sum statement.
\end{proof}

The absolute value in \eqref{eq:coeff} is an integer: $n!/
\prod a_i!(i!)^{a_i}$ counts partitions of an $n$-element labelled set
with $a_i$ blocks of size $i$, and $(m-1)!$ is an integer.
The substitution $x_i=z$ gives
\begin{equation}\label{eq:stirling}
 L_n(z,z,\ldots)=\sum_{m=1}^n(-1)^{m-1}(m-1)!
 \left\{\begin{matrix}n\\m\end{matrix}\right\}z^m.
\end{equation}
Indeed, summing the set-partition counts above over exponent vectors
with $m$ blocks gives the Stirling number in braces. This is the
signed A028246 specialization mentioned in A263634.

\begin{corollary}\label{cor:appell}
Set $x_0=1$ and define $P_n(z)$ by
$F(t)e^{zt}=\sum_{n\geq0}P_n(z)t^n/n!$. Then
\begin{align}
 x_{n+1}&=\sum_{j=0}^n\binom nj L_{j+1}x_{n-j},\label{eq:moment}\\
 P_{n+1}(z)&=zP_n(z)+\sum_{j=0}^n\binom nj L_{j+1}P_{n-j}(z).
 \label{eq:appell}
\end{align}
\end{corollary}
\begin{proof}
Multiplying $(\log F)'$ by $F$ gives $F'$, so coefficient extraction
gives \eqref{eq:moment}. Differentiating $F(t)e^{zt}$ with respect
to $t$ gives $(z+(\log F)')F(t)e^{zt}$, which gives
\eqref{eq:appell}. These identities include $n=0$.
\end{proof}
The following additional source identities have short formal proofs.
Write $I_m=m![t^m]F^{-1}$. Then
\begin{equation}\label{eq:inverse}
 R(n,k)=\sum_{r=0}^{n-1}\binom{n-1}{r}x_{k+r}I_{n-1-r}.
\end{equation}
This follows by multiplying $F^{(k)}$ and $F^{-1}$. The $I_m$ are
the reciprocal polynomials of A133314. If $U_n(z)$ has generating
function $F^{-1}e^{zt}$, the proof of \eqref{eq:appell} gives
$U_{n+1}=zU_n-\sum_{j=0}^n\binom nj L_{j+1}U_{n-j}$.
For any $G(0)=1$, let $T_G(z^n)=n![t^n]G(t)e^{zt}$ and extend
linearly. Applying this to the generating function shows
$T_GT_H=T_{GH}$, hence $T_FT_{F^{-1}}$ is the identity. This proves
the source's umbral inverse claim.

For the three exponential families in the source, define
\[
 \begin{split}
 \sum_{n\geq0}C_n(a)\frac{t^n}{n!}&=\exp\!\left(\sum_{i\geq1}a_i\frac{t^i}{i}\right),\\
 \sum_{n\geq0}B_n(b)\frac{t^n}{n!}&=\exp\!\left(\sum_{i\geq1}b_i\frac{t^i}{i!}\right),\\
 \sum_{n\geq0}H_n(c)\frac{t^n}{n!}&=\exp\!\left(\sum_{i\geq1}c_it^i\right).
 \end{split}
\]
These are the cycle-index, Bell, and refined Lah families of
A036039, A036040, and A130561, respectively~\cite{oeisfamilies}.
Taking formal logarithms gives
\begin{equation}\label{eq:families}
 L_n(C_1,\ldots,C_n)=(n-1)!a_n,\quad
 L_n(B_1,\ldots,B_n)=b_n,\quad
 L_n(H_1,\ldots,H_n)=n!c_n.
\end{equation}
Conversely, $\exp(\log F)=F$ gives the inverse substitutions
\[
 x_n=C_n\left(\frac{L_1}{0!},\ldots,\frac{L_n}{(n-1)!}\right)
 =B_n(L_1,\ldots,L_n)
 =H_n\left(\frac{L_1}{1!},\ldots,\frac{L_n}{n!}\right).
\]
Thus the source's six substitutions A--F follow, with their stated
factorial conventions. With the Faber convention
$-\log(1+\sum_{i\geq1}b_it^i)=\sum_{n\geq1}\Phi_n(b)t^n/n$
of A263916, coefficient comparison gives
$L_n=-(n-1)!\Phi_n(x_1/1!,\ldots,x_n/n!)$.

The matrix formula in the source also follows from \eqref{eq:appell}.
Put $S_{n,k}=\binom nk x_{n-k}$ for $0\leq k\leq n$, and put
$M_{n,k}=1$ when $k=n+1$ and $0$ otherwise. Since
$P_n=\sum_{k=0}^nS_{n,k}z^k$,
\[
 D^jP_n=\frac{n!}{(n-j)!}P_{n-j},\qquad
 \mathcal R=z+\sum_{j\geq0}\frac{L_{j+1}}{j!}D^j
\]
give $\mathcal RP_n=P_{n+1}$ and $S\mathcal R=MS$ in the monomial
basis. Here $S^{-1}_{n,k}=\binom nk I_{n-k}$ by $FF^{-1}=1$.
Therefore the matrix is $S^{-1}MS$, with first superdiagonal $1$
and entries $\binom nk L_{n+1-k}$ for $k\leq n$.
All rows in these products have finite support.

Finally, let $T_n(w)$ denote the row polynomial of A130850, with
\[
 \sum_{n\geq0}T_n(w)\frac{t^n}{n!}
 =\frac{w}{(1+w)e^{-wt}-1}.
\]
On substituting $x_i=z^{i-1}$, the logarithmic derivative is
$z/(1+(z-1)e^{-zt})$. Replacing $z$ by $-w$ identifies it with
the preceding generating function at $-t$. Hence
$L_{n+1}(1,z,z^2,\ldots)=(-1)^nT_n(-z)$, the signed specialization
in the source's 14 November 2016 comment.
The rational calculations may be performed in $\Q(z)[[t]]$;
the resulting identities are polynomial identities in $z$.

\section{Normalization and repeated Appell zeros}
The 17 May 2020 comment in A263634 considers
\begin{equation}\label{eq:canonical}
 F(t)=e^{bt}\prod_{j=1}^d(1-t/\rho_j)e^{t/\rho_j},
 \qquad \rho_j\in\mathbb R\setminus\{0\}.
\end{equation}
Expanding each logarithm gives
\begin{equation}\label{eq:normalization}
 L_1=b,\qquad L_n=-(n-1)!\sum_{j=1}^d\rho_j^{-n}\quad(n\geq2).
\end{equation}
Thus, for the power sums $c_n=\sum_j\rho_j^{-n}$ used in that
comment, the correct relation is $c_n=-L_n/(n-1)!$, rather than
$c_n=-L_n$. The same relation holds for an infinite product with
$\sum_j\rho_j^{-2}<\infty$: for sufficiently small $|t|$, the sum
of the absolute values of the logarithmic tails is bounded by a
constant times $|t|^2\sum_j\rho_j^{-2}$. This justifies coefficient
comparison. A Gaussian factor $e^{-at^2}$ changes only $L_2$, by
$-2a$.

The hypotheses in that comment also permit repeated Appell zeros.
In \eqref{eq:canonical} take $d=1$, $\rho_1=1$, and $b=0$, so that
$F=(1-t)e^t$. Then
\[
 L_3=-2,\quad c_3=1,\qquad
 P_n(z)=(z+1)^n-n(z+1)^{n-1},\qquad
 P_3(z)=(z+1)^2(z-2).
\]
This gives both a normalization counterexample and a double Appell
zero under the stated assumptions. More generally, if
$F(t)=e^{\beta t}p(t)$ with $p(0)=1$ and $\deg p=d$, then
\begin{equation}\label{eq:multiplicity}
 P_n(z)=p(D)(z+\beta)^n
\end{equation}
by coefficient extraction. For $n>d$, the zero at $z=-\beta$
has multiplicity exactly $n-d$: the highest derivative
term is a nonzero multiple of $(z+\beta)^{n-d}$ and all other
terms have larger exponent. This is the exponential-polynomial
exception in Craven and Csordas~\cite[Proposition 4.1(2)]{craven}
when $F$ belongs to the Laguerre--P\'olya class and $\beta\ne0$;
the displayed algebraic argument also covers arbitrary $p$ and $\beta=0$.
In particular universal simplicity requires additional hypotheses.

\section{Source coverage}
Dates below refer to A263634 unless an entry is named explicitly.
\begin{center}\small
\begin{tabular}{@{}p{.48\linewidth}p{.45\linewidth}@{}}
\toprule
Source statement & Coverage\\\midrule
Kurkov, FORMULA, 30 March 2025 & Alternative proof: Theorem~\ref{thm:main}; all-column Leonov--Shiryaev specialization.\\
Row count and row sums; examples & Proved: Corollary~\ref{cor:coeff}.\\
Constant-variable specialization, 6 November 2015 & Reproved: \eqref{eq:stirling}; signed A028246.\\
Moment/Appell identities, 13 September 2016 & Reproved; prior formulas: Novak (2012), Copeland (2015).\\
Reciprocal and matrix comments, 2015--2016 & Proved above: $T_F^{-1}=T_{F^{-1}}$ and $S^{-1}MS$.\\
Six substitutions A--F, 8--12 September 2016 & Proved: \eqref{eq:families} and inverse substitutions.\\
A127671, A263916: cumulant/Faber relation & Reproved: coefficient comparison above.\\
Signed A130850, 14 November 2016 & Proved by the displayed generating functions.\\
Power-sum normalization, 17 May 2020 & Refuted-and-corrected: \eqref{eq:normalization}.\\
Universal simplicity in the same comment & Refuted: $(z+1)^2(z-2)$; the finite-product multiplicity is \eqref{eq:multiplicity}, as in Craven--Csordas (1994).\\
Stable graphs, Hopf algebras, and historical attributions & Not covered as separate results.\\
A263633, A178867; Lagrange inversion in A134264 & Not covered as separate results.\\
\bottomrule
\end{tabular}
\end{center}
The remaining cross-references A000041, A007318, A129185, A133314,
and A154921 supply, respectively, partition counts, Pascal coefficients,
the shift matrix, reciprocal polynomials, and the specialization
$F=2-e^t$ of the reciprocal Appell family.

\section{Mathematical verification coverage}
The pinned file
\href{https://github.com/the-omega-institute/trureturing/blob/6a3b29c88128eeed17f731783f03e7f9ff92ae62/D5/S3/Combinatorics/PowerSeries/KurkovA263634LogPolynomials.lean}{\small\texttt{KurkovA263634LogPolynomials.lean}}
defines \texttt{sourceSeries}, \texttt{IsSourceFamily}, and the array
\texttt{R}. Its declaration \texttt{result} states existence of the logarithmic family and
$R(n,1)=L_n$ for every $n\geq1$ over $\Q[x_1,x_2,\ldots]$.
The helper \texttt{quotient\_jet} expresses the all-column identity
\eqref{eq:quotient}. Its logarithmic-derivative dependency is
\href{https://github.com/the-omega-institute/trureturing/blob/6a3b29c88128eeed17f731783f03e7f9ff92ae62/D5/S3/Combinatorics/PowerSeries/PositiveNewtonSums.lean}{\texttt{PositiveNewtonSums.lean}}.
The general substitution, coefficient, Appell, and zero-multiplicity
statements have the written proofs given above. The accompanying
\texttt{verify.py} uses exact rational polynomial arithmetic to compare
the quadratic recurrence, formal logarithm, partition coefficients,
and quotient coefficients for $1\leq n\leq10$, $1\leq k\leq6$;
it also checks the moment and Appell identities through degree $9$
and the first $44$ coefficients of A263634. These are finite checks.

\begin{thebibliography}{99}
\bibitem{oeis} The OEIS Foundation Inc., \emph{A263634: logarithmic
polynomials}, revision \#166, 5 November 2025. Conjecture by Mikhail
Kurkov, FORMULA, 30 March 2025.
\url{https://oeis.org/A263634}.
\bibitem{oeis127671} The OEIS Foundation Inc., \emph{A127671:
cumulant expansion numbers}, revision \#88, 4 February 2026.
\url{https://oeis.org/A127671}.
\bibitem{xpaul} xpaul, answer to \emph{Simplify recurrence
$f'_{n-1}=f_n-f_{n-1}f_1$}, Mathematics Stack Exchange,
21 March 2019. \url{https://math.stackexchange.com/a/3156943}.
\bibitem{iv} IV\_, answer on repeated derivatives for
$f'_{k-1}=f_k-f_{k-1}f_1$, Mathematics Stack Exchange,
15 July 2020. \url{https://math.stackexchange.com/a/3758094}.
\bibitem{novak} J. Novak, with illustrations by M. LaCroix,
\emph{Three lectures on free probability}, arXiv:1205.2097v1 (2012),
Section 1.2, pp.~4--6. \url{https://arxiv.org/abs/1205.2097v1}.
\bibitem{copeland} T. Copeland, \emph{The creation / raising
operators for Appell sequences}, 2015, Part II, p.~3.
\url{https://tcjpn.wordpress.com/wp-content/uploads/2015/12/appell-raising-op1.pdf}.
\bibitem{craven} T. Craven and G. Csordas, \emph{Differential
operators of infinite order and the distribution of zeros of entire
functions}, J. Math. Anal. Appl. \textbf{186} (1994), 799--820,
Proposition 4.1(2). \url{https://doi.org/10.1006/jmaa.1994.1334}.
\bibitem{oeisfamilies} The OEIS Foundation Inc., entries
\href{https://oeis.org/A036039}{A036039},
\href{https://oeis.org/A036040}{A036040},
\href{https://oeis.org/A130561}{A130561},
\href{https://oeis.org/A133314}{A133314},
\href{https://oeis.org/A028246}{A028246},
\href{https://oeis.org/A130850}{A130850}, and
\href{https://oeis.org/A263916}{A263916}, accessed 14 November 2026.
\bibitem{ls} V. P. Leonov and A. N. Shiryaev, \emph{On a method
of calculation of semi-invariants}, Theory Probab. Appl. \textbf{4}
(1959), 319--329. \url{https://doi.org/10.1137/1104031}.
\bibitem{btv} T. H. Beuman, A. M. Turner and V. Vitelli,
\emph{Extrema statistics in the dynamics of a non-Gaussian random field},
arXiv:1211.0993v1 (2012), Appendix A, (A10)--(A15);
Phys. Rev. E \textbf{87} (2013), 022142.
\url{https://arxiv.org/abs/1211.0993v1}.
\bibitem{burchardt} A. Burchardt, \emph{Algebras with two
multiplications and their cumulants}, arXiv:1803.09322v1 (2018),
Section 1.2 and Theorem 1.9.
\url{https://arxiv.org/abs/1803.09322v1}.
\bibitem{anshelevich} M. Anshelevich, \emph{Appell polynomials
and their relatives}, arXiv:math/0311043v3 (2004), Section 1, (1), and Section 2.2.
\url{https://arxiv.org/abs/math/0311043v3}.
\bibitem{mo} user231922, \emph{Similar recursions for complicated
partition polynomials}, MathOverflow question 489828,
21 March 2025. \url{https://mathoverflow.net/q/489828}.
\end{thebibliography}
\medskip\noindent\small
\copyright\ 2026 Haobo Ma and Wenlin Zhang.
Licensed under Creative Commons Attribution 4.0 International
(CC BY 4.0).
\end{document}
